\documentclass[11pt,a4paper]{article}
\usepackage[T1]{fontenc}
\usepackage[utf8]{inputenc}
\usepackage{lmodern,amsmath,amssymb}
\usepackage[margin=25mm]{geometry}
\usepackage[hidelinks]{hyperref}
\hypersetup{pdftitle={The blocking criterion is monotonically worse in the block size: blocking gains nothing wi},pdfauthor={navstokgap research working notes}}
\newcommand{\tightlist}{\setlength{\itemsep}{0pt}\setlength{\parskip}{0pt}}
\setlength{\emergencystretch}{3em}
\setcounter{secnumdepth}{0}
\begin{document}
\hypertarget{the-blocking-criterion-is-monotonically-worse-in-the-block-size-blocking-gains-nothing-without-a-change-of-coupling}{%
\section{The blocking criterion is monotonically worse in the block
size: blocking gains nothing without a change of
coupling}\label{the-blocking-criterion-is-monotonically-worse-in-the-block-size-blocking-gains-nothing-without-a-change-of-coupling}}

Reading Yarotsky's proof (Commun. Math. Phys. 261 (2006) 799, §2,
passage level) settles the input left open by
\href{blocking-step-obstruction.md}{the blocking note}: the proof is a
polymer expansion, its Lemma 1 bounds the non-classical evolution of an
excited region \(I\) by \((2\alpha e^{t_0\beta/\alpha})^{|I|}\) and its
Lemma 3 gives configuration weights exponentially damped in the size of
the excited region, so the connected structure that the blocking step
needed is present, and the criterion applies to a lattice of blocks with
the block Hamiltonians as the classical part. The conclusion is negative
and sharp. With blocks of \(M^3\) sites the criterion reads
\[\beta_{\rm block}=\frac{\text{straddling norm per block}}{\text{block gap}}
=\frac{24M^2N}{g^2\,\delta(g;M)}\ \le\ \beta_*,\] and \(\delta(g;M)\)
degrades or stays flat as \(M\) grows: at strong coupling the block gap
is the single flux-loop energy \(2C_2g^2\), independent of \(M\), so
\(\beta_{\rm block}\propto M^2\); in the small-volume regime
\(\delta\simeq\delta_1g^{2/3}/M\), so \(\beta_{\rm block}\propto M^3\).
Against the direct criterion \(\beta_{\rm direct}=32N/(C_2g^4)\) of
\href{strong-coupling-uniform-gap.md}{the T2 note} the ratio is
\(\beta_{\rm block}/\beta_{\rm direct}\simeq\tfrac38M^2\), already above
one at \(M=2\). \textbf{Real-space blocking with a criterion of this
form is strictly worse than no blocking at all}, at every coupling, and
iterating it cannot reach weaker coupling. This corrects Sections 4 and
5 of \href{blocking-step-obstruction.md}{the blocking note}, which
assumed that the connected estimate would buy a fixed threshold and an
induction: the estimate is available, and it buys nothing, because the
boundary grows like \(M^2\) while the gap it is measured against does
not grow at all. Any gain must come from recognizing the blocked
Hamiltonian as a theory of the same form with a \textbf{larger effective
coupling}, which is the renormalization step itself. Constants explicit;
nothing promoted.

\hypertarget{the-connected-structure-is-present-in-the-proof}{%
\subsection{1. The connected structure is present in the
proof}\label{the-connected-structure-is-present-in-the-proof}}

Yarotsky's Theorem 1 is proved by writing
\[e^{-t_0H_\Lambda}=\sum_{I\subset\Lambda}T_{\Lambda,I},\qquad
T_{\Lambda,I}=\sum_{J\subset I}(-1)^{|I|-|J|}\,e^{-t_0(H_{\Lambda,0}+\sum_{x\in J}\phi_x)},\]
so that \(T_{\Lambda,I}\) collects the contributions in which every site
of \(I\) is touched by the perturbation. Two estimates carry the
argument (§2, passage level).

\begin{itemize}
\tightlist
\item
  \textbf{Lemma 1.} \(\|T_I\|\le(2\alpha e^{t_0\beta/\alpha})^{|I|}\),
  proved by analytic continuation of the couplings
  \(\phi_x\to z_x\phi_x\), the Hille--Yosida bound
  \(\|e^{-t_0H_J(z_J)}\|\le e^{t_0|J|\beta/\alpha}\) on the numerical
  range, and a multidimensional Schwarz lemma using that \(T_I(z_I)\)
  vanishes whenever any \(z_x=0\).
\item
  \textbf{Lemma 3.} The weight of a space-time configuration
  \(C=\{(I_k,J_k)\}\) obeys
  \(|w(C)|\le\prod_k(2\alpha e^{t_0\beta/\alpha})^{|I_k|}e^{-t_0(|J_k|-|\Lambda_0|^3|I_k|)}\),
  the classically excited sites contributing \(e^{-t_0}\) each through
  the local gap.
\end{itemize}

Both are exponential in the size of the excited region, which is exactly
the polymer structure that makes the expansion sum over
\textbf{connected} clusters. So the hypothetical estimate of
\href{blocking-step-obstruction.md}{the blocking note} §3, replacing
\((\sum_p\|w_p\|)^2\) by \(\sum_p\|w_p\|^2\), is precisely what the
expansion produces, and the way to use it is to apply the theorem itself
to the blocked system in place of a re-derived Schur bound.

\hypertarget{the-criterion-on-a-lattice-of-blocks}{%
\subsection{2. The criterion on a lattice of
blocks}\label{the-criterion-on-a-lattice-of-blocks}}

Identify the blocks of \(M^3\) sites with the sites of a coarse cubic
lattice, as in \href{blocking-step-obstruction.md}{the blocking note}
§1. Take as classical part the normalized block Hamiltonians
\[h_\alpha=\frac{H_\alpha-E_0^\alpha}{\Delta_M},\qquad
\Delta_M=\frac{\hbar c}{a}\,\delta(g;M),\] each with non-degenerate
ground state \(\Omega_\alpha\) (T1) and unit gap, and as perturbation
the straddling plaquettes,
\[\phi_\alpha=\frac1{\Delta_M}\sum_{p\ \rm straddling\ at\ \alpha}w_p,
\qquad
\|\phi_\alpha\|\le\frac{6M^2\cdot4N\hbar c/(ag^2)}{\Delta_M}
=\frac{24M^2N}{g^2\,\delta(g;M)}=:\beta_{\rm block}.\] The perturbation
is bounded, so \(\alpha=0\) in Yarotsky's condition (2) and the
criterion is \(\beta_{\rm block}\le\beta_*(3,\{0,1\}^3)\), with the same
constants as in \href{strong-coupling-uniform-gap.md}{the T2 note}. The
classical part is a single-block operator, hence a function of its own
spectral partition of unity, and the partition contains the projection
onto \(\Omega_\alpha\); the interaction range is one coarse lattice
spacing.

\hypertarget{the-criterion-degrades-with-m}{%
\subsection{\texorpdfstring{3. The criterion degrades with
\(M\)}{3. The criterion degrades with M}}\label{the-criterion-degrades-with-m}}

Two regimes fix the behaviour of \(\delta(g;M)\).

\emph{Strong coupling.} The lowest gauge-invariant excitation of a block
is a single plaquette flux loop, of energy
\(2C_2(R_{\min})g^2\hbar c/a\) independent of the block size, as in
\href{mass-gap-obligations-lattice.md}{the obligations map} §3. So
\(\delta(g;M)\to2C_2g^2\) and
\[\beta_{\rm block}\simeq\frac{24M^2N}{2C_2g^4}=\frac{12M^2N}{C_2g^4},\]
\[\frac{\beta_{\rm block}}{\beta_{\rm direct}}
\simeq\frac{12M^2N/(C_2g^4)}{32N/(C_2g^4)}=\frac{3M^2}{8},\] using
\(\beta_{\rm direct}=32N/(C_2g^4)\) from the link-level check of
\href{strong-coupling-uniform-gap.md}{the T2 note}. At \(M=2\) the ratio
is \(3/2\), and it grows quadratically thereafter.

\emph{Small volume.} For \(g\) small and a block of side \(Ma\) the gap
is the zero-mode gap of \href{low-dimensional-mass-gap.md}{G07}
Proposition 7, \(\Delta_M=\delta_1g^{2/3}\hbar c/(Ma)\), that is
\(\delta(g;M)=\delta_1g^{2/3}/M\). Then
\[\beta_{\rm block}\simeq\frac{24M^3N}{\delta_1\,g^{8/3}},\] growing
like \(M^3\) and diverging as \(g\to0\).

In both regimes the numerator grows like the block surface, \(M^2\),
while the denominator stays flat or shrinks. The criterion therefore has
its best value at \(M=1\), which is the direct application already used
for T2.

\textbf{Proposition.} For every \(M\ge2\) and every \(g>0\),
\(\beta_{\rm block}(M)\ge\beta_{\rm block}(1)\) whenever
\(\delta(g;M)\le M^2\delta(g;1)\), which holds in both regimes above and
follows in general from the upper bounds of
\href{polyakov-average-gap-bound.md}{the Polyakov note}, where the gap
of a box of side \(L\) is at most of order \(\hbar c/L\) times a
coupling factor. Blocking therefore never improves the criterion.

\hypertarget{what-this-rules-out-and-what-it-leaves}{%
\subsection{4. What this rules out, and what it
leaves}\label{what-this-rules-out-and-what-it-leaves}}

The negative result is specific: \textbf{a criterion that compares the
total inter-block coupling with the block gap cannot be improved by
blocking}, because coupling scales with area and the gap does not scale
with volume. It rules out the naive induction of
\href{blocking-step-obstruction.md}{the blocking note} §5, in which one
hoped that a fixed threshold, once reached, would propagate to all
larger scales.

What survives is the observation that made the induction attractive: the
effective coupling of an asymptotically free theory grows toward the
infrared. To use it, the blocked Hamiltonian must be \textbf{recognized
as a theory of the same type with a new coupling}, so that the criterion
is re-applied with \(g_{n+1}>g_n\) and with the block again playing the
role of a single site of unit spacing. That identification is the
renormalization-group step; it requires a change of variables at each
scale, not a projection. Yarotsky's own Theorem 3 does precisely this in
its setting: he treats the AKLT model by showing that \emph{on a large
length scale it is a relatively bounded perturbation of a classical
model}, which works because the frustration-free structure makes the
block ground state exactly known. The Kogut--Susskind Hamiltonian at
intermediate coupling has no such structure, and supplying a substitute
is the open problem.

\hypertarget{correction-to-the-previous-note}{%
\subsection{5. Correction to the previous
note}\label{correction-to-the-previous-note}}

\href{blocking-step-obstruction.md}{The blocking note} §§4--5 concluded
that a connected estimate would give a one-step inequality valid above a
fixed threshold \(g^4\gtrsim8\sqrt{6C}MN/(\gamma C_2)\) and that the
induction would then close. Sections 1--3 above show that the connected
estimate is available and that the resulting threshold is worse than the
direct one by the factor \(3M^2/8\), so the induction does not close by
this route. The finite-step counting of §5 of that note,
\(n\simeq(2b_0\log2)^{-1}(g_{\rm UV}^{-2}-g_{\rm thr}^{-2})\), retains
its meaning as a count of renormalization-group steps, with the proviso
that each step must redefine the coupling rather than merely project.

\hypertarget{consequence-for-state}{%
\subsection{6. Consequence for STATE}\label{consequence-for-state}}

The blocking route in its projection form is closed, with the reason
stated quantitatively: boundary grows like \(M^2\), gap does not grow.
The route survives only in its renormalization form, where each step
redefines the theory, which is the constructive problem. The nearest
remaining tractable questions are on the other side of the ledger: the
stability theorem without frustration-freeness named in
\href{lieb-robinson-kogut-susskind.md}{the Lieb--Robinson note} §5, and
the fixed-lattice small-volume theorem (1a) named in
\href{schur-error-ultraviolet.md}{the Schur note} §5.
\end{document}
